\documentclass[11pt]{article}
\usepackage[margin=1.00in]{geometry}
\usepackage[T1]{fontenc}
\usepackage{lmodern,amsmath,amssymb,amsthm,mathtools,microtype,booktabs}
\usepackage{xurl}\usepackage[hidelinks]{hyperref}
\usepackage{enumitem}
\setlist{nosep}
\newtheorem{theorem}{Theorem}[section]
\newtheorem{lemma}[theorem]{Lemma}
\newtheorem{corollary}[theorem]{Corollary}
\theoremstyle{remark}\newtheorem{remark}[theorem]{Remark}
\newcommand{\R}{\mathbb R}\newcommand{\C}{\mathbb C}\newcommand{\N}{\mathbb N}
\newcommand{\ii}{\mathrm i}\newcommand{\e}{\mathrm e}\newcommand{\ph}{\varphi}
\newcommand{\norm}[1]{\left\lVert #1\right\rVert}\newcommand{\abs}[1]{\left|#1\right|}
\DeclareMathOperator{\ord}{ord}\DeclareMathOperator{\Rea}{Re}\DeclareMathOperator{\meas}{meas}
\title{\textbf{Spectral coercivity of multiscale\\smoothed prime-counting errors}\\[5pt]\large A consequence of Dirichlet GRH without zero-independence assumptions}
\author{Research note prepared for independent review}
\date{26 November 2026}
\begin{document}
\maketitle
\begin{abstract}
For a fixed modulus and any finite family of distinct multiplicative scales, we prove that the centered, square-root-normalized errors of a smoothed prime-power count have strictly positive definite joint logarithmic covariance. This remains true after deleting an arbitrary finite set of zero frequencies. More quantitatively, in every sufficiently high positive-ordinate band $(T,T+H]$, for $K\log(2T)\le H\le T$, the covariance quadratic form is bounded below by $cH/(T^4\log(2T))$ times the squared Euclidean norm of the coefficients and above by $CH\log^2(2T)/T^4$ times that norm. The proof uses Dirichlet GRH, the twice-denominated explicit formula and fixed-argument Landau sums. A prime in each residue class can be chosen to isolate each scale without interference from prime powers at other scales. No simplicity, disjointness, or rational linear independence of zero ordinates is required. Consequences include uniform two-sided oscillation in every direction with positive lower logarithmic density, and a full-dimensional convex hull of limiting error vectors. The formulation and proof are independently developed here; priority in the literature is not certified.
\end{abstract}

\section{Connection with the verified GRH development}
The source is the readable edition of \texttt{lsernau/grh-lean} \cite{repo}, a shorter rewriting of the Alp\"oge--Claude proof. Its \texttt{src/Statements.lean} imports Mathlib and states the zero assertions independently of the proof infrastructure. In \texttt{src/Main.lean}, \texttt{New.grh} yields the half-plane theorem \path{LFunction_ne_zero_of_half_lt_re}. Reflection yields the zero-on-line theorem \path{LFunction_eq_zero_re}, for every Dirichlet character, primitive or not. The zeta case is also stated as Mathlib's \texttt{RiemannHypothesis}. Only the classical Dirichlet version of GRH is used here, not GRH for arbitrary automorphic or Dedekind $L$-functions.

The repository's \texttt{GUIDE.md} and \texttt{New/Reduction/MainInequality.lean} explain the comparison of an upper bound for the covolume and index of a family of short integral forms with a lower bound from twisted theta lifts. Cancellation of a main term leads to \texttt{star\_inequality}, in the source's notation $m\,\mathrm{rem}\le fa(\log b)^5$. The subsequent file \texttt{New/Reduction/OneSidedBound.lean} controls the remainder of a prime-power sum with kernel $(1/n-1/x)$ and weight $\Lambda(n)(1-(n/d))$ in a residue class. Finally \texttt{New/Reduction/Zeros/Criterion.lean} uses a one-sided Landau criterion and quadratic twists to obtain the zero-free half-plane. Its Mellin residues have denominator $s(1-s)$, precisely the denominator that will become positive on the critical line below. The companion \texttt{PrimeNumberTheorem.lean} expresses the smoothed remainder as the integral of a counting error.

Our use of these ideas reverses the direction: given the zero location, we ask what oscillations of the same smoothing are \emph{forced} and whether several scales or residue classes can cancel them. We replace the special quadratic weight by arbitrary residue-class weights and exploit a second, classical Landau formula (the sum of $x^\rho$ over a vertical band, not the one-sided criterion). Wikipedia's treatment of prime distribution, analytic criteria, and related zero-correlation questions \cite{wiki} provides further motivation. The new argument does not modify or reprove the upstream formal proof.

\paragraph{Proof status.} The repository reports compilation and axiom checks, as well as independent proof-checker acceptance for the five statements. The Lean FRO/Mathlib statement and Rubinstein's verification record give independent evidence \cite{lean,waterloo}. This note uses those conclusions as a theorem input and classical analytic number theory as described below. We inspected the guide, README, theorem statements and the relevant reduction files, not every proof module; neither the upstream build nor a Lean formalization of this new result was run here. The upstream proof still lacks a complete human-read exposition. For a reader who wishes to bracket that status, every assertion in this note can equally be read as a theorem conditional only on classical Dirichlet GRH.

\section{Definitions and statement}
Fix an integer $q\ge1$, write $G=(\mathbb Z/q\mathbb Z)^\times$, $d=\ph(q)$, and let $\widehat G$ denote its characters. For $q=1$ use the trivial group, whose unique character is identically one. For $a\in G$ and $x\ge1$ set
\begin{align*}
 S_a(x)&=\sum_{\substack{n\le x\\ n\equiv a\pmod q}}\Lambda(n)\left(\frac1n-\frac1x\right),\\
 c_a&=\lim_{x\to\infty}\left(S_a(x)-\frac{\log x}{d}\right),\\
 R_a(x)&=S_a(x)-\frac{\log x}{d}-c_a .
\end{align*}
The limit exists, as proved in Lemma~\ref{lem:exp}. Here $\Lambda(p^k)=\log p$ for primes $p$ and integers $k\ge1$, and $\Lambda(n)=0$ otherwise. In particular the count includes prime powers, not just primes.

Write $\chi^*$ for the primitive character inducing $\chi$, with conductor dividing $q$. Let $m_\chi(\gamma)$ be the multiplicity of $1/2+\ii\gamma$ as a nontrivial zero of $L(s,\chi^*)$ (zero if absent). Let $\Gamma$ be the union of the ordinates of these zeros, as a \emph{set}, so common ordinates occur only once. For any real $\gamma$ define
\[
 v_a(\gamma)=\frac1d\sum_{\chi\in\widehat G}\overline{\chi(a)}m_\chi(\gamma),\qquad b_a=4v_a(0).
\]
The conjugation identity $m_{\bar\chi}(-\gamma)=m_\chi(\gamma)$ implies $v_a(-\gamma)=\overline{v_a(\gamma)}$; in particular $b_a\in\R$. Define the centered signal on logarithmic time by
\[
 Z_a(t)=\e^{t/2}R_a(\e^t)-b_a\qquad(t\ge0).
\]
Fix $r\ge1$ distinct real numbers $h_1,\ldots,h_r$ (the multiplicative scales are $\e^{h_j}$). Let $D=rd$. For an array $w=(w_{j,a})\in\C^D$ put
\[
 A_w(\gamma)=\sum_{j=1}^r\sum_{a\in G} w_{j,a}\e^{\ii\gamma h_j}v_a(\gamma).
\]
For real $w$, the positive-frequency band $B_{T,H}=\{\gamma\in\Gamma:T<\gamma\le T+H\}$ has associated real signal
\[
 F_{w;T,H}(t)=2\Rea\sum_{\gamma\in B_{T,H}}
 \frac{A_w(\gamma)}{1/4+\gamma^2}\e^{\ii\gamma t}.
\]
By integrating the finitely many exponentials, its mean square on $[0,Y]$ as $Y\to\infty$ is
\begin{equation}\label{eq:energy}
 V_{T,H}(w)=2\sum_{\gamma\in B_{T,H}}\frac{|A_w(\gamma)|^2}{(1/4+\gamma^2)^2}.
\end{equation}
Use the right side to define the same quadratic form for complex $w$ as well.

\begin{theorem}[Multiscale spectral coercivity]\label{thm:main}
Assume Dirichlet GRH. For fixed $q$ and distinct $h_1,\dots,h_r$, there exist positive constants $c,C,K,T_0$, depending only on $q$ and the shifts, such that for every real $T\ge T_0$, every $K\log(2T)\le H\le T$, and every $w\in\C^D$,
\begin{equation}\label{eq:main}
 \boxed{\quad \frac{cH}{T^4\log(2T)}\norm{w}_2^2\ \le V_{T,H}(w)\ \le
 \frac{CH\log^2(2T)}{T^4}\norm{w}_2^2.\quad}
\end{equation}
In particular every such band has at least $D$ distinct positive zero ordinates, and the vectors $(\e^{\ii\gamma h_j}v_a(\gamma))_{j,a}$ for those ordinates span $\C^D$.
\end{theorem}
The short-window threshold is logarithmic in height for fixed modulus and scales; no uniformity in $q$, in the number of scales, or as two shifts approach equality is asserted. The statement is meaningful also for $q=1$ and arbitrarily many distinct scales. The exact covariance formula retains common zeros with their actual multiplicities, rather than counting their contributions as independent.

\section{Analytic inputs and the explicit formula}
We use the classical analytic continuation and functional equation of primitive Dirichlet $L$-functions, their Euler products, and the local zero count $N_\chi(U+1)-N_\chi(U-1)\ll_q\log(2+|U|)$, where the count is with multiplicity. These are standard (e.g. \cite{mv,banks}); the same estimates hold for $\zeta$. In particular $\sum_\gamma m_\chi(\gamma)/(1/4+\gamma^2)<\infty$. GRH provides the exact real part $1/2$ of all nontrivial zeros. All sums over zeros here include central zeros unless explicitly excluded.

\begin{lemma}[Absolutely convergent smoothed explicit formula]\label{lem:exp}
For $t\ge\log 2$, uniformly over $a\in G$,
\begin{equation}\label{eq:explicit}
 Z_a(t)=\sum_{\substack{\gamma\in\Gamma\\\gamma\ne0}}
       \frac{v_a(\gamma)}{1/4+\gamma^2}\e^{\ii\gamma t}
       +O_q((1+t)\e^{-t/2}).
\end{equation}
The series converges absolutely and uniformly for all real $t$.
\end{lemma}
\begin{proof}
First take a primitive $\chi$ and denote $D_\chi(s)=-L'(s,\chi)/L(s,\chi)$. Set $\delta_\chi=1$ for the conductor-one character and $0$ otherwise. The elementary Mellin inversion for the ramp $\max(1-1/y,0)$ gives, on the line $\Rea z=1$,
\begin{equation}\label{eq:perron}
 \sum_{n\le x}\Lambda(n)\chi(n)\left(\frac1n-\frac1x\right)
 =\frac1{2\pi\ii}\int_{1-\ii\infty}^{1+\ii\infty}
        D_\chi(1+z)\frac{x^z}{z(z+1)}\,dz.
\end{equation}
Indeed the Dirichlet series is absolutely convergent there and the kernel is integrable. At $y=1$ the ramp is zero, so no endpoint convention is needed. Shift the line to $\Rea z=-3/2$. The functional equation and Stirling's estimate give $D_\chi(-1/2+\ii u)\ll_q\log(2+|u|)$; this line contains no zeros or poles. The local zero count permits horizontal contours at heights tending to infinity (positive and negative separately) whose distance from every zero ordinate is $\gg_q1/\log(2+|u|)$. The standard local logarithmic derivative estimate (\cite[Lemmas 12.1, 12.6]{mv}) on $-1/2\le\Rea(1+z)\le2$ then gives $O_q(\log^2(2+|u|))$ for $D_\chi$. Thus for fixed $x$ the horizontal integrals vanish, and the new vertical integral is $O_q(x^{-3/2})$ uniformly for $x\ge2$.

At $z=0$, write $D_\chi(1+z)=\delta_\chi/z+k_\chi+O_\chi(z)$. The residue is $\delta_\chi\log x+k_\chi-\delta_\chi$. At $z=\rho-1$ for a nontrivial zero of multiplicity $m$, the residue is
\[
 -\frac{m x^{\rho-1}}{(\rho-1)\rho}
 =\frac{m x^{-1/2+\ii\gamma}}{1/4+\gamma^2},\qquad \rho=1/2+\ii\gamma.
\]
The remaining possible pole crossed is $z=-1$: the kernel pole and, for a primitive even nonprincipal character, a trivial zero at $1+z=0$. Its order is at most two; the residue is $O_\chi((1+\log x)/x)$. This proves the formula with $\delta_\chi\log x+c_\chi$ as main term and $c_\chi=k_\chi-\delta_\chi$. The zero series is absolutely convergent by the count above.

For imprimitive $\chi\pmod q$, its value differs from $\chi^*$ only at integers with a prime factor dividing $q$; in the Mangoldt sum only prime powers at those primes matter. The discarded contribution from each $p\mid q$ is a constant (the absolutely convergent sum $\log p\sum_{k\ge1}\chi^*(p)^k/p^k$) plus $O_q((1+\log x)/x)$, by the geometric tail and the at most $\log x/\log p$ ramp terms. Thus only the constant and error change, and the nontrivial zeros are still those of $\chi^*$. Averaging with $d^{-1}\overline{\chi(a)}$ by character orthogonality yields $S_a(x)=d^{-1}\log x+c_a+x^{-1/2}\sum_\gamma v_a(\gamma)x^{\ii\gamma}/(1/4+\gamma^2)+O_q((1+\log x)/x)$. This also proves existence and reality of $c_a$. Multiplication by $x^{1/2}$ and removal of the central term gives \eqref{eq:explicit}.
\end{proof}

\begin{lemma}[Fixed-argument Landau formula]\label{lem:landau}
For a fixed real $u\ge2$, uniformly for $1<T<T+H\le 2T$, one has
\[
 \sum_{T<\gamma\le T+H}m_\chi(\gamma)\e^{\ii\gamma\log u}
 =-\frac{H}{2\pi\sqrt u}\Lambda(u)\chi^*(u)+O_{q,u}(\log(2T)),
\]
where the product on the right is interpreted as zero unless $u$ is an integral prime power.
\end{lemma}
\begin{proof}
Banks's extension of the Landau--Gonek formula \cite[Theorem 1.2]{banks}, applied to each primitive character, gives the sum of $u^\rho$ over the open ordinate band with main term $-H\Lambda(u)\chi^*(u)/(2\pi)$. For fixed $u$ the distance to the nearest \emph{different} prime power is positive and all the terms of the error are $O_{q,u}(\log(2T))$. For conductor one the same estimate is the classical fixed-$u$ Landau formula (difference of two prefix sums). The choice of half-open instead of open endpoints changes the result by $O_{q,u}(\log(2T))$ by the local zero count. Divide by $\sqrt u$, using GRH. No independence or simplicity hypothesis is used.
\end{proof}

\section{Isolating shifts by primes}
\begin{lemma}[Prime probes]\label{lem:probes}
For each $a\in G$ one can choose a prime $p_a\equiv a\pmod q$ such that for every pair $j,k$, the real number
\[
 u_{a,j,k}=p_a\e^{h_k-h_j}\ge2
\]
is an integral prime power if and only if $j=k$ (when it is $p_a$).
\end{lemma}
\begin{proof}
For the finitely many off-diagonal ratios $\e^{h_k-h_j}$ that are rational, express them as reduced positive fractions $n/m\ne1$. Exclude the finitely many primes dividing any numerator or denominator. Also require $p_a\min_{j,k}\e^{h_k-h_j}\ge2$. Dirichlet's theorem provides arbitrarily large primes in each reduced class; choose one avoiding all exclusions. An irrational ratio times an integer prime is irrational. For a rational ratio with $m>1$, $p_a n/m$ is not an integer, since $p_a$ is coprime to $m$. If $m=1$, then $n\ge2$, and $p_a n$ has at least two distinct prime factors, since $p_a$ does not divide $n$. Thus all off-diagonal main terms vanish. The diagonal is the chosen prime.
\end{proof}

\begin{theorem}[Recovery of the coefficient array]\label{thm:recovery}
With the primes in Lemma~\ref{lem:probes}, for every $w\in\C^D$,
\begin{equation}\label{eq:recovery}
 \sum_{\gamma\in B_{T,H}} A_w(\gamma)\e^{\ii\gamma(\log p_a-h_j)}
 =-H\frac{\log p_a}{2\pi\sqrt{p_a}}w_{j,a}
      +O_{q,h,p}(\norm{w}_2\log(2T)).
\end{equation}
The error is uniform in $j,a,w$ and $1<T<T+H\le2T$.
\end{theorem}
\begin{proof}
Expand $A_w$ into its $r d^2$ terms. For fixed $j,a$, the left side equals
\[
 \frac1d\sum_{k,b,\chi} w_{k,b}\overline{\chi(b)}
 \sum_{\gamma\in B_{T,H}}m_\chi(\gamma)\e^{\ii\gamma\log u_{a,j,k}}.
\]
Use Lemma~\ref{lem:landau} for each fixed argument $u_{a,j,k}$. By Lemma~\ref{lem:probes}, only $k=j$ can contribute a main term. Then $u=p_a$ is coprime to $q$, so $\chi^*(p_a)=\chi(a)$. Orthogonality gives $d^{-1}\sum_\chi \overline{\chi(b)}\chi(a)=\mathbf 1_{a=b}$, isolating $w_{j,a}$. The sum of all error bounds is at most a fixed constant times $\sum_{k,b}|w_{k,b}|\log(2T)\le\sqrt D\norm{w}_2\log(2T)$.
\end{proof}
Thus any array is recoverable, with error $O_{q,h,p}(\log(2T)/H)\norm{w}_2$ per entry, from a band of aggregated zero data. In particular the recovery converges uniformly on bounded coefficient sets when $H/\log T\to\infty$.

\begin{proof}[Proof of Theorem~\ref{thm:main}]
Put $L=\log(2T)$, and $\alpha=\min_a(\log p_a)/(2\pi\sqrt{p_a})>0$. Choose an entry with $|w_{j,a}|\ge\norm{w}_2/\sqrt D$. If the error in \eqref{eq:recovery} is bounded by $E L\norm{w}_2$, take $K\ge \max(1,2E\sqrt D/\alpha)$. Then for $H\ge KL$ the absolute value of the left side of \eqref{eq:recovery} is at least $H\alpha\norm{w}_2/(2\sqrt D)$. Set $M_\gamma=\sum_\chi m_\chi(\gamma)$. The unit-interval zero count gives constants $C_0,C_1>0$ depending only on $q$ such that for $H\ge1$,
\[
 |B_{T,H}|\le\sum_{\gamma\in B_{T,H}}M_\gamma\le C_0 HL,
 \qquad \max_{\gamma\in B_{T,H}}M_\gamma\le C_1L.
\]
Cauchy--Schwarz on the \emph{distinct} ordinates in the probe sum gives
\begin{equation}\label{eq:lower}
 \sum_{\gamma\in B_{T,H}}|A_w(\gamma)|^2\ge
 \frac{\alpha^2}{4DC_0}\frac{H}{L}\norm{w}_2^2.
\end{equation}
Since $1/4+\gamma^2\le5T^2$ for $\gamma\le2T$ and $T\ge1$, the lower bound in \eqref{eq:main} holds with $c=\alpha^2/(50 DC_0)$.

For the upper bound, $|v_a(\gamma)|\le M_\gamma/d$, so $|A_w(\gamma)|\le(\sqrt D/d)M_\gamma\norm{w}_2$. Hence
\[
 \sum_{\gamma\in B_{T,H}}|A_w(\gamma)|^2\le
 \frac{D}{d^2}\norm{w}_2^2 C_1 L\,C_0 HL.
\]
The denominator in \eqref{eq:energy} is at least $T^4$, proving the upper bound with $C=2DC_0 C_1/d^2$. For $T$ large enough the allowed range of $H$ is nonempty. The lower bound says that the complex-linear map $w\mapsto(A_w(\gamma))_{\gamma\in B_{T,H}}$ is injective, proving the span and cardinality assertions. Neither Cauchy--Schwarz nor character orthogonality required the zeros to be distinct across characters.
\end{proof}

\paragraph{A determinant corollary.} For real arrays, let $Q_{T,H}$ denote the real symmetric matrix representing \eqref{eq:energy}. All its eigenvalues lie between the two bounds of \eqref{eq:main} for a unit vector. Thus
\[
 \det Q_{T,H}\ge\left(\frac{cH}{T^4 L}\right)^D,
 \qquad \mathrm{cond}_2(Q_{T,H})\le (C/c)L^3 .
\]
The same eigenvalue estimates apply to the Hermitian form on complex arrays (with the conventional conjugation of its coefficient vector). This finite-dimensional determinant estimate is another parallel with the upstream comparison-of-covolumes idea, but its proof here is purely spectral.

\section{Logarithmic covariance and genuine oscillation}
Let $J$ be any finite subset of the positive ordinates in $\Gamma$ (it may be empty). Define a real $D$-vector $U_J(t)$ for all sufficiently large $t$ by its entries
\[
 (U_J(t))_{j,a}=Z_a(t+h_j)-2\Rea\sum_{\gamma\in J}
      \frac{v_a(\gamma)\e^{\ii\gamma(t+h_j)}}{1/4+\gamma^2}.
\]
To avoid irrelevant domain restrictions when integrating from zero, extend it by zero where some $t+h_j<0$; changing a bounded initial interval cannot affect any limit. Set $z(\gamma)=(\e^{\ii\gamma h_j}v_a(\gamma))_{j,a}$ and $d_\gamma=1/4+\gamma^2$. Lemma~\ref{lem:exp} gives $U_J(t)=F_J(t)+o(1)$ in Euclidean norm, where
\[
 F_J(t)=2\Rea\sum_{\substack{\gamma\in\Gamma,\gamma>0\\\gamma\notin J}}
                  z(\gamma)\e^{\ii\gamma t}/d_\gamma.
\]
This is a bounded uniform limit of trigonometric polynomials, hence uniformly almost periodic. For real $w$ the mean of $w\cdot F_J$ is zero. By termwise integration for finite sums and uniform convergence, its mean square is exactly $2\sum_{\gamma>0,\gamma\notin J}|A_w(\gamma)|^2/d_\gamma^2$. The same limits hold for $U_J$ because the difference is bounded eventually and tends to zero. In particular the logarithmic covariance matrix exists and is
\begin{equation}\label{eq:cov}
 Q_J=\lim_{Y\to\infty}\frac1Y\int_0^Y U_J(t)U_J(t)^{\mathsf T}dt
 =2\Rea\sum_{\substack{\gamma>0\\\gamma\notin J}}
             \frac{z(\gamma)\overline{z(\gamma)}^{\mathsf T}}{d_\gamma^2}.
\end{equation}
Its smallest eigenvalue $\lambda_J$ is strictly positive: take any allowed band above every ordinate in $J$, and apply Theorem~\ref{thm:main} to that band's contribution. This proof also works for deletion of any bounded set of ordinates, since the union of zeros is locally finite.

\begin{corollary}[Uniform two-sided density]\label{cor:density}
There exist $\eta,\delta>0$, depending only on $q,h,J$, such that for every real unit array $w$, each of the inequalities $w\cdot U_J(\log x)\ge\eta$ and $w\cdot U_J(\log x)\le-\eta$ holds on a set of $x\ge1$ of lower logarithmic density at least $\delta$. In particular no nonzero linear contrast becomes $o(1)$ or eventually one-sided after centering and deleting these frequencies.
\end{corollary}
\begin{proof}
Let $M=2\sum_{\gamma>0}\norm{z(\gamma)}_2/d_\gamma$, finite and positive. Write $f(t)=w\cdot F_J(t)$, so $|f|\le M$ and its mean is zero. The inequality $f^2\le M|f|$ gives $\liminf_{Y\to\infty}Y^{-1}\int_0^Y |f|\,dt\ge\lambda_J/M$. Since $f_+=(|f|+f)/2$, each positive and negative part has lower mean at least $\lambda_J/(2M)$. With $\theta=\lambda_J/(4M)$, pointwise $f_+\le\theta+M\mathbf1_{f\ge\theta}$, and the same argument applies to $-f$. Thus both sign sets have lower time density at least $\lambda_J/(4M^2)$. As $\norm{U_J-F_J}_2=o(1)$, take $\eta=\lambda_J/(8M)$ and $\delta=\lambda_J/(4M^2)$ for the original signal. The substitution $t=\log x$ identifies time density with $\liminf_{X\to\infty}(\log X)^{-1}\int_1^X\mathbf1_E(x)\,dx/x$.
\end{proof}

\begin{corollary}[Full-dimensional convex hull]\label{cor:hull}
The closed convex hull of the set of limit points of $U_J(t)$ as $t\to\infty$ contains the closed Euclidean ball of radius $\eta$ centered at zero.
\end{corollary}
\begin{proof}
The cluster set is nonempty and compact since the vector is bounded at large times. For each unit $w$, Corollary~\ref{cor:density} gives arbitrarily large $t$ with $w\cdot U_J(t)\ge\eta$; a convergent subsequence yields a point with the same inequality. Thus the support function of the convex hull is at least $\eta$ in every unit direction. The separation theorem for closed convex sets implies the ball inclusion.
\end{proof}

\begin{remark}[A version requiring no central-zero information]
If $\sum_j w_{j,a}=0$ for each $a$, the bias $b_a$ cancels identically even if unknown. With $J=\varnothing$, the nonzero array then gives two-sided oscillations of the direct linear combination $\sum_{j,a} w_{j,a}\sqrt{\e^{h_j}x}\,R_a(\e^{h_j}x)$. For instance for any residue class and $c>1$, the difference $\sqrt{cx}\,R_a(cx)-\sqrt{x}\,R_a(x)$ exceeds a positive constant and falls below its negative on sets of positive lower logarithmic density. There is no need to evaluate any zeros in order to define this difference.
\end{remark}

\section{Examples, novelty and limitations}
\paragraph{Probe example.} For $q=12$ and scales $1,2,3/2,5$ the reduced classes are $1,5,7,11$. Primes can be chosen with $p\ge10$ excluding $2,3,5$, namely $13,17,19,11$ respectively. All $p$ times off-diagonal scale ratios are either nonintegers or composite numbers with at least two different prime factors. The accompanying script checks the diagonal/off-diagonal condition exactly using rational arithmetic. For these 16 coordinates, Theorem~\ref{thm:main} guarantees rank 16 over $\C$ in every band of the stated sufficiently large sizes.

\paragraph{Zeta-only check.} For $q=1$, the constant is $c=-1-\gamma_{\rm E}$, where $\gamma_{\rm E}$ is Euler's constant, and there is no central zero. The supplementary script evaluates a truncated zeta series using the first 60 positive zeros and compares it with the exact finite prime-power expression (up to floating-point roundoff) at a few large $x$. It also computes illustrative finite covariance matrices for shifts $0,\log 2,\log 3$. For $x=10^3,3\cdot10^3,10^4,3\cdot10^4,10^5,3\cdot10^5,10^6$, the discrepancy between the direct value and the truncated series plus the $z=-1$ correction $\log(2\pi)/\sqrt x$ is less than $5\cdot 10^{-4}$ in every tested case. For the selected ordinates 41--60 the three finite covariance eigenvalues are approximately $7.43\cdot10^{-8}$, $1.00\cdot10^{-7}$ and $1.21\cdot10^{-7}$. These are numerical sanity checks of signs, normalization, and the Gram construction, not a certification of zero completeness or a proof of the large-height bounds.

\paragraph{Prior work and novelty.} Classical prime-race work uses explicit formulas and limiting logarithmic distributions; see Rubinstein--Sarnak \cite{rs}, Martin--Ng \cite{mn}, and Devin \cite{devin}. Banks \cite{banks} uses Landau sums to distinguish the zero multisets of two characters in short windows, and earlier work of Fujii and Raghunathan also distinguishes such zero multisets (see Banks's introduction). The recent Banks--Loftus preprint \cite{bl} addresses a different nonvanishing question for $L$-functions evaluated at other zeros. We make no novelty claim for the individual classical tools. The proposed new contribution is their combination into a uniform coercivity estimate for \emph{all} real or complex coefficient arrays across arbitrary finite distinct logarithmic shifts, with explicit arithmetic isolation of each shift and class, and its covariance/deletion consequences. A targeted literature search found related results but did not establish priority; independent review is required.

\paragraph{What is not proved.} The result does not imply that each unsmoothed prime race ordering occurs, that the cluster set itself contains a ball (only its convex hull does), that the zeros are simple or rationally independent, or that the constants are small. It concerns a smoothed prime-power count with its mean removed. The $O$-constants can in principle be extracted from classical estimates for fixed data; no numerical $T_0,K,c,C$ are certified here. The proof is classical, human-readable mathematics, not a kernel-checked Lean development.

\section{Further research questions}
\begin{enumerate}
\item Can the gap of three powers of $\log T$ between the eigenvalue bounds be closed using information about pair correlations and multiplicities? With stronger hypotheses, what are the sharp eigenvalue asymptotics for arbitrary shifts?
\item Can one bound $K$ uniformly as $q$ or the number of scales grows, or in terms of the minimum gap between shifts? What are efficient, near-optimal probe primes?
\item Can the convex-hull result be strengthened to interior of the actual limiting support, or to inclusivity for suitably centered smoothed races, without a zero-ordinate independence conjecture? Covariance alone cannot establish this.
\item Can the method handle compactly supported logarithmic convolution kernels, or families of Artin and automorphic L-functions with distinguishable prime coefficients and suitable zero-location theorems?
\item Formalize the finite prime-probe and Gram arguments in Lean; then bridge to the explicit formula, Landau--Gonek estimate, and the upstream five theorems. The short-window Landau theorem and quantitative local zero count are the major analytic interfaces to verify.
\end{enumerate}

\begin{thebibliography}{99}\small
\bibitem{repo} L. Sernau et al., \emph{GRH for Dirichlet L-functions, in Lean 4}, readable and compact editions, November 2026. \url{https://github.com/lsernau/grh-lean}. See also the source audit in the archive.
\bibitem{wiki} Wikipedia contributors, \emph{Riemann hypothesis}, accessed 26 November 2026. \url{https://en.wikipedia.org/wiki/Riemann_hypothesis}.
\bibitem{lean} Lean FRO and Mathlib Initiative, \emph{Lean verification of the Riemann Hypothesis}, November 2026. \url{https://leanprover-community.github.io/blog/posts/riemann-hypothesis/}.
\bibitem{waterloo} M. Rubinstein, \emph{Alp\"oge/Claude proof of the Generalized Riemann Hypothesis}, verification record, accessed 26 November 2026. \url{https://www.math.uwaterloo.ca/~mrubinst/alpoge_claude_grh/}.
\bibitem{mv} H. L. Montgomery and R. C. Vaughan, \emph{Multiplicative Number Theory I: Classical Theory}, Cambridge University Press, 2007, especially chapters 10--12.
\bibitem{banks} W. D. Banks, \emph{Comparing zeros of distinct Dirichlet L-functions}, Publ. Math. Debrecen 106 (2025), 327--335. arXiv:2302.07073, version accessed 26 November 2026; Theorem 1.2. \url{https://arxiv.org/abs/2302.07073}.
\bibitem{rs} M. Rubinstein and P. Sarnak, \emph{Chebyshev's bias}, Experimental Mathematics 3 (1994), 173--197.
\bibitem{mn} G. Martin and N. Ng, \emph{Inclusive prime number races}, arXiv:1710.00088. \url{https://arxiv.org/abs/1710.00088}.
\bibitem{devin} L. Devin, \emph{Limiting properties of the distribution of primes in an arbitrarily large number of residue classes}, Canadian Mathematical Bulletin, DOI:10.4153/S0008439520000089.
\bibitem{bl} W. Banks and K. Loftus, \emph{Unshared zeros of Dirichlet L-functions}, arXiv:2607.22930 (2026). \url{https://arxiv.org/abs/2607.22930}.
\end{thebibliography}
\end{document}
